\subsection{Nonlinear Operator Theory Methods}\label{sec:nonlinear-operator-theory-methods}

In this section we give a number of results on nonlinear mappings from a Banach space $X$ to its dual:
\begin{equation}\label{eq:nonlinear-banach-to-dual-mapping}
X\ni v\mapsto T(v)\in X^*.
\end{equation}
We shall usually use the term \emph{mapping} to refer to a nonlinear mapping and reserve the term \emph{operator} for when we are assuming a mapping is linear. However, the reader should be warned that this is an affectation adopted for this book which is not used throughout the literature. Authors use the term ``operator\textquotesingle\textquotesingle\ for both linear and nonlinear mappings, the assumption that an operator is linear is often left unspecified in contexts where only linear operators are studied. (We do this in Chapter~\ref{chap:operator-theory}.)

\subsubsection{Mappings on Finite-Dimensional Spaces}\label{subsec:finite-dimensional-nonlinear-mappings}

In this section we study mappings $\bm{f}:\mathbb{R}^n\to\mathbb{R}^n$. Our goal is a better understanding of the problem of ``$n$ equations in $n$ unknowns\textquotesingle\textquotesingle:
\begin{equation}\label{eq:finite-dimensional-nonlinear-system}
\bm{f}(\bm{u})=\bm{a}.
\end{equation}
We already have discussed a ``local\textquotesingle\textquotesingle\ result for the problem, the inverse function theorem, in Chapter~\ref{chap:introduction}; but we desire something stronger. As an example of what we expect, consider the following result when $n=1$.

\begin{Theorem}
\label{thm:scalar-continuous-monotone-surjectivity}
Suppose $f:\mathbb{R}\to\mathbb{R}$ is continuous and satisfies
\begin{equation}\label{eq:scalar-continuous-surjectivity-growth-condition}
f(u)\to\pm\infty\quad\text{as }u\to\pm\infty.
\end{equation}
Then, for every $a\in\mathbb{R}$ the equation
\begin{equation}\label{eq:scalar-continuous-surjectivity-equation}
f(u)=a
\end{equation}
has a solution $u\in\mathbb{R}$. Furthermore, if $f$ is monotone increasing, then, for each $a$, the set of solutions forms a closed interval; if $f$ is strictly monotone, the solution is unique.
\end{Theorem}

The proof is an elementary application of the intermediate value theorem and the definition of monotone and continuous functions.

In order to get an analogue of this result for mappings from $\mathbb{R}^n$ to $\mathbb{R}^n$ we must generalize both the growth condition \eqref{eq:scalar-continuous-surjectivity-growth-condition} and the definition of monotone increasing (order preserving).

\begin{Definition}
\label{def:finite-dimensional-coercive-mapping}
We say that a function $\bm{f}:\mathbb{R}^n\to\mathbb{R}^n$ is \textbf{coercive} if
\begin{equation}\label{eq:finite-dimensional-nonlinear-coercivity-definition}
\dfrac{\bm{f}(\bm{u})\cdot\bm{u}}{|\bm{u}|}\to\infty\quad\text{as }|\bm{u}|\to\infty.
\end{equation}

\end{Definition}

\begin{Definition}
\label{def:finite-dimensional-monotone-mapping}
We say that a function $\bm{f}:\mathbb{R}^n\to\mathbb{R}^n$ is monotone if
\begin{equation}\label{eq:finite-dimensional-monotonicity-definition}
(\bm{f}(\bm{u})-\bm{f}(\bm{v}))\cdot(\bm{u}-\bm{v})\geq0\quad\text{for all }\bm{u},\bm{v}\in\mathbb{R}^n.
\end{equation}
We say that $\bm{f}$ is strictly monotone if the inequality in \eqref{eq:finite-dimensional-monotonicity-definition} is strict whenever $\bm{u}\neq\bm{v}$.
\end{Definition}

We first study the implications of monotonicity. We begin with the following immediate consequence of the definition.

\begin{Theorem}
\label{thm:finite-dimensional-strict-monotonicity-uniqueness}
If $\bm{f}$ is strictly monotone, then the equation
\begin{equation}\label{eq:finite-dimensional-strictly-monotone-equation}
\bm{f}(\bm{u})=\bm{a}
\end{equation}
has at most one solution.
\end{Theorem}

To get a result when $\bm{f}$ is only monotone we first prove the following lemma.

\begin{Lemma}
\label{lem:finite-dimensional-minty-equivalence}
Let $\bm{a}\in\mathbb{R}^n$ be given and let $\bm{f}:\mathbb{R}^n\to\mathbb{R}^n$ be continuous. Then if $\bm{u}\in\mathbb{R}^n$ is a solution of the variational inequality
\begin{equation}\label{eq:finite-dimensional-minty-variational-inequality}
(\bm{f}(\bm{v})-\bm{a})\cdot(\bm{v}-\bm{u})\geq0\quad\text{for every }\bm{v}\in\mathbb{R}^n,
\end{equation}
it follows that $\bm{u}$ is also a solution of equation \eqref{eq:finite-dimensional-strictly-monotone-equation}.

In addition, if $\bm{f}$ is monotone, then every solution of equation \eqref{eq:finite-dimensional-strictly-monotone-equation} is also a solution of the variational inequality \eqref{eq:finite-dimensional-minty-variational-inequality}.
\end{Lemma}

\begingroup
\renewcommand{\endProof}{\par\nobreak\hfill$\blacksquare$\par\vspace{0.1cm}}
\begin{Proof}
Suppose $\bm{f}$ is continuous and \eqref{eq:finite-dimensional-minty-variational-inequality} holds. Then, for any $\bm{w}\in\mathbb{R}^n$, let $\bm{v}=\bm{u}+t\bm{w}$ for $t>0$. This gives us (after dividing \eqref{eq:finite-dimensional-minty-variational-inequality} by $t$)
\begin{equation}\label{eq:finite-dimensional-minty-directional-inequality}
(\bm{f}(\bm{u}+t\bm{w})-\bm{a})\cdot\bm{w}\geq0\quad\text{for all }\bm{w}\in\mathbb{R}^n.
\end{equation}
Letting $t\to0$ gives us
\begin{equation}\label{eq:finite-dimensional-minty-directional-limit}
(\bm{f}(\bm{u})-\bm{a})\cdot\bm{w}\geq0\quad\text{for all }\bm{w}\in\mathbb{R}^n.
\end{equation}
Using $\bm{w}=\pm(\bm{f}(\bm{u})-\bm{a})$ gives us $\bm{f}(\bm{u})=\bm{a}$.

Now, suppose $\bm{f}$ is monotone and \eqref{eq:finite-dimensional-strictly-monotone-equation} holds. Then
\begin{equation}\label{eq:finite-dimensional-minty-monotonicity-converse}
(\bm{f}(\bm{v})-\bm{a})\cdot(\bm{v}-\bm{u})=(\bm{f}(\bm{v})-\bm{f}(\bm{u}))\cdot(\bm{v}-\bm{u})\geq0\quad\text{for every }\bm{v}\in\mathbb{R}^n,
\end{equation}
by the definition of monotonicity.
\end{Proof}
\endgroup

As an immediate consequence of this we get the following.

\begin{Theorem}
\label{thm:finite-dimensional-monotone-solution-set-closed-convex}
Let $\bm{a}\in\mathbb{R}^n$ and let $\bm{f}$ be continuous and monotone. Then the set $\mathcal{K}$ of solutions of $\bm{f}(\bm{u})=\bm{a}$ is closed and convex.
\end{Theorem}

\begingroup
\renewcommand{\endProof}{\par\nobreak\hfill$\blacksquare$\par\vspace{0.1cm}}
\begin{Proof}
For any fixed $\bm{v}\in\mathbb{R}^n$ let
\begin{equation}\label{eq:monotone-solution-set-halfspaces}
S_{\bm{v}}=\{\bm{u}\in\mathbb{R}^n\mid(\bm{f}(\bm{v})-\bm{a})\cdot(\bm{v}-\bm{u})\geq0\}.
\end{equation}
Note that each set $S_{\bm{v}}$ is a closed half-space (and is hence convex). By our Lemma~\ref{lem:finite-dimensional-minty-equivalence} we can write the solution set as
\begin{equation}\label{eq:monotone-solution-set-halfspace-intersection}
\mathcal{K}=\bigcap_{\bm{v}\in\mathbb{R}^n}S_{\bm{v}}.
\end{equation}
Since the arbitrary intersection of closed sets is closed and the arbitrary intersection of convex sets is convex, our theorem is proved.
\end{Proof}
\endgroup

We now examine the consequences of coercivity.

\begin{Theorem}
\label{thm:finite-dimensional-coercive-surjectivity}
Let $\bm{f}:\mathbb{R}^n\to\mathbb{R}^n$ be continuous and coercive. Then for every $\bm{a}\in\mathbb{R}^n$, equation \eqref{eq:finite-dimensional-strictly-monotone-equation} has a solution $\bm{u}\in\mathbb{R}^n$.
\end{Theorem}

In order to prove this we need the following important result.

\begin{Theorem}{Brouwer fixed point theorem}
\label{thm:brouwer-fixed-point}
Let $C$ be a compact, convex, nonempty subset of $\mathbb{R}^n$, and suppose $\bm{f}$ is a continuous function that maps $C$ into $C$. Then $\bm{f}$ has a fixed point in $C$; i.e., there exists $\bm{u}\in C$ such that
\begin{equation}\label{eq:brouwer-fixed-point-equation}
\bm{f}(\bm{u})=\bm{u}.
\end{equation}

\end{Theorem}

\begingroup
\renewcommand{\endProof}{\par\nobreak\hfill$\blacksquare$\par\vspace{0.1cm}}
\begin{Proof}
For the general case, we refer the reader to the literature (cf. \cite{DS}), but when $n=1$ the proof is easy to see graphically. In this case the hypotheses of the theorem state that we have a continuous function $f$ which maps an interval $[a,b]$ into itself. The graph of the function must lie in the box $[a,b]\times[a,b]$ in the $(x,y)$-plane. There are only three possibilities: $f(a)=a$, $f(b)=b$ or the graph of $f$ starts to the left of the line $x=y$ and ends on the right. Since the function is continuous the graph must cross the line $x=y$ somewhere. Any crossing gives us a fixed point.
\end{Proof}
\endgroup

We now prove Theorem~\ref{thm:finite-dimensional-coercive-surjectivity}.

\begingroup
\renewcommand{\endProof}{\par\nobreak\hfill$\blacksquare$\par\vspace{0.1cm}}
\begin{Proof}
We begin by noting that it is sufficient to show that $\bm{f}(\bm{u})=0$ has a solution. To see this let an arbitrary $\bm{a}\in\mathbb{R}^n$ be given, and define
\begin{equation}\label{eq:coercive-surjectivity-shifted-mapping}
\bm{f}_{\bm{a}}(\bm{u}):=\bm{f}(\bm{u})-\bm{a}.
\end{equation}
We now note that $\bm{f}_{\bm{a}}$ is coercive if and only if $\bm{f}$ is, since\[\begin{aligned}
\dfrac{\bm{f}_{\bm{a}}(\bm{u})\cdot\bm{u}}{|\bm{u}|}&=\dfrac{\bm{f}(\bm{u})\cdot\bm{u}}{|\bm{u}|}-\dfrac{\bm{a}\cdot\bm{u}}{|\bm{u}|}\\
&\geq\dfrac{\bm{f}(\bm{u})\cdot\bm{u}}{|\bm{u}|}-\dfrac{|\bm{a}||\bm{u}|}{|\bm{u}|}\\
&=\dfrac{\bm{f}(\bm{u})\cdot\bm{u}}{|\bm{u}|}-|\bm{a}|.
\end{aligned}\]Thus, if we can show that for every coercive function, $\bm{f}(\bm{u})=0$ has a solution, it immediately follows that $\bm{f}(\bm{u})=\bm{a}$ has one for every $\bm{a}\in\mathbb{R}^n$.

We now convert $\bm{f}(\bm{u})=0$ to a fixed point problem. As a first step, we define $\bm{g}:\mathbb{R}^n\to\mathbb{R}^n$ by
\begin{equation}\label{eq:coercive-surjectivity-fixed-point-mapping}
\bm{g}(\bm{u}):=\bm{u}-\bm{f}(\bm{u}).
\end{equation}
Now note that
\begin{equation}\label{eq:coercive-surjectivity-fixed-point-pairing}
\bm{g}(\bm{u})\cdot\bm{u}=|\bm{u}|^2-\bm{f}(\bm{u})\cdot\bm{u}.
\end{equation}
Since $\bm{f}$ is coercive, there exists $R>0$ such that
\begin{equation}\label{eq:coercive-surjectivity-large-radius-strict-bound}
\bm{g}(\bm{u})\cdot\bm{u}<|\bm{u}|^2\quad\text{whenever }|\bm{u}|\geq R.
\end{equation}
We now define
\begin{equation}\label{eq:coercive-surjectivity-ball-retraction}
\bm{r}(\bm{v}):=\begin{cases}\bm{v},&|\bm{v}|\leq R,\\\dfrac{R\bm{v}}{|\bm{v}|},&|\bm{v}|>R.\end{cases}
\end{equation}
Finally, we let
\begin{equation}\label{eq:coercive-surjectivity-retracted-fixed-point-mapping}
\bm{h}(\bm{u}):=\bm{r}(\bm{g}(\bm{u})).
\end{equation}
We note that $\bm{h}$, as the composition of continuous functions, is continuous. Furthermore, if we let
\begin{equation}\label{eq:coercive-surjectivity-open-ball}
B:=\{\bm{u}\in\mathbb{R}^n\mid|\bm{u}|<R\},
\end{equation}
then $\bm{h}$ maps the closed ball $\overline B$ to itself. Thus, by the Brouwer fixed point theorem, $\bm{h}$ has a fixed point $\bm{u}\in\overline B$.

If $\bm{u}\in\partial B$, then
\begin{equation}\label{eq:coercive-surjectivity-boundary-fixed-point-norms}
R=|\bm{u}|=|\bm{h}(\bm{u})|=|\bm{r}(\bm{g}(\bm{u}))|.
\end{equation}
Thus, by the definition of $\bm{r}$, we must have
\begin{equation}\label{eq:coercive-surjectivity-boundary-radius-ratio}
\dfrac{R}{|\bm{g}(\bm{u})|}\leq1.
\end{equation}
But we now use \eqref{eq:coercive-surjectivity-large-radius-strict-bound} to get
\begin{equation}\label{eq:coercive-surjectivity-boundary-contradiction}
|\bm{u}|^2=\bm{u}\cdot\bm{u}=\bm{h}(\bm{u})\cdot\bm{u}=\dfrac{R}{|\bm{g}(\bm{u})|}\bm{g}(\bm{u})\cdot\bm{u}<\dfrac{R}{|\bm{g}(\bm{u})|}|\bm{u}|^2\leq|\bm{u}|^2.
\end{equation}
This is a contradiction. Thus, we must have
\begin{equation}\label{eq:coercive-surjectivity-interior-fixed-point-norms}
R>|\bm{u}|=|\bm{h}(\bm{u})|=|\bm{r}(\bm{g}(\bm{u}))|.
\end{equation}
Hence, $\bm{r}(\bm{g}(\bm{u}))=\bm{g}(\bm{u})$ and so
\begin{equation}\label{eq:coercive-surjectivity-interior-retraction-identity}
\bm{u}=\bm{h}(\bm{u})=\bm{r}(\bm{g}(\bm{u}))=\bm{g}(\bm{u})=\bm{u}-\bm{f}(\bm{u}).
\end{equation}
Finally, this gives us
\begin{equation}\label{eq:coercive-surjectivity-zero-equation-conclusion}
\bm{f}(\bm{u})=0,
\end{equation}
and completes the proof.
\end{Proof}
\endgroup

Combining the results of this section gives us

\begin{Corollary}
\label{cor:finite-dimensional-coercive-strict-monotone-bijection}
If $\bm{f}:\mathbb{R}^n\to\mathbb{R}^n$ is continuous, strictly monotone and coercive, then for every $\bm{a}\in\mathbb{R}^n$ there exists a unique $\bm{u}\in\mathbb{R}^n$ such that\[\bm{f}(\bm{u})=\bm{a}.\]
\end{Corollary}

\subsubsection{Monotone Mappings on Banach Spaces}\label{subsec:monotone-banach-mappings}

In this section we will see that we can expand some of the ideas of the previous section to infinite-dimensional spaces. We will assume $X$ is a real, reflexive Banach space and we study mappings $T:X\to X^*$. For an element $g\in X^*$ we will use the standard inner product notation and write $(g,v)$ for $g(v)$. \emph{In the remaining sections of this chapter, we will, in order to make the text easier to follow, give proofs only in the case where $X$ and $X^*$ are separable.}

We begin by giving a string of definitions which generalize concepts from functions on finite-dimensional spaces to mappings on reflexive Banach spaces.

\begin{Definition}
\label{def:banach-nonlinear-bounded-continuous-mapping}
We say that a mapping $T:X\to X^*$ is \textbf{bounded} if it maps bounded sets in $X$ to bounded sets in $X^*$. The mapping is \textbf{continuous} if for every $u\in X$ we have
\begin{equation}\label{eq:banach-nonlinear-mapping-continuity-definition}
\|T(u)-T(v)\|_{X^*}\to0\quad\text{whenever }\|u-v\|_X\to0.
\end{equation}

\end{Definition}

\begin{Definition}
\label{def:banach-nonlinear-monotone-mapping}
We say that a mapping $T:X\to X^*$ is monotone if
\begin{equation}\label{eq:banach-nonlinear-mapping-monotonicity-definition}
(T(u)-T(v),u-v)\geq0\quad\text{for all }u,v\in X,
\end{equation}
and strictly monotone if this inequality is strict whenever $u\neq v$.
\end{Definition}

\begin{Definition}
\label{def:banach-nonlinear-coercive-mapping}
We say that a mapping $T:X\to X^*$ is coercive if
\begin{equation}\label{eq:banach-nonlinear-mapping-coercivity-definition}
\dfrac{(T(u),u)}{\|u\|}\to\infty\quad\text{as }\|u\|\to\infty.
\end{equation}

\end{Definition}

\begin{Remark}
\label{rem:finite-banach-pairing-generalization}
Using the definitions above, the results of the previous section can easily be extended to mappings $T$ from an $n$-dimensional Banach space $X$ to another $n$-dimensional Banach space $Y$. The only real trick here is in defining the appropriate bilinear form
\begin{equation}\label{eq:finite-banach-spaces-appropriate-bilinear-pairing}
X\times Y\ni x,y\mapsto\langle x,y\rangle\in\mathbb{R}
\end{equation}
to replace the ``inner product\textquotesingle\textquotesingle
\begin{equation}\label{eq:banach-space-dual-inner-product-notation}
X\times X^*\ni x,y\mapsto(x,y)\in\mathbb{R}
\end{equation}
in the definitions in this section and the dot product in the previous section. The details are left to the reader.
\end{Remark}

The proofs of Lemma~\ref{lem:finite-dimensional-minty-equivalence} and Theorem~\ref{thm:finite-dimensional-monotone-solution-set-closed-convex} were based on monotonicity and were in no way dependent on the dimension of the spaces involved. Thus, we can get the following two analogous results using only minor changes of notation in the previous proofs.

\begin{Lemma}
\label{lem:banach-minty-equivalence}
Let $g\in X^*$ be given and let $T:X\to X^*$ be continuous. Then if $u\in X$ is a solution of the variational inequality
\begin{equation}\label{eq:banach-minty-variational-inequality}
(T(v)-g,v-u)\geq0\quad\text{for every }v\in X,
\end{equation}
then $u$ satisfies the equation
\begin{equation}\label{eq:banach-minty-equation}
T(u)=g.
\end{equation}
In addition, if $T$ is monotone, then every solution of the equation \eqref{eq:banach-minty-equation} is also a solution of the variational inequality \eqref{eq:banach-minty-variational-inequality}.
\end{Lemma}

\begin{Theorem}
\label{thm:banach-monotone-solution-set-closed-convex}
Let $g\in X^*$ and let $T:X\to X^*$ be continuous and monotone. Then the set $\mathcal{K}$ of solutions of $T(u)=g$ is closed and convex.
\end{Theorem}

Our first infinite-dimensional analogue of Theorem~\ref{thm:finite-dimensional-coercive-surjectivity} is the following.

\begin{Theorem}{Browder-Minty}
\label{thm:browder-minty-surjectivity}
Let $X$ be a real, reflexive Banach space and let $T:X\to X^*$ be bounded, continuous, coercive and monotone. Then for any $g\in X^*$ there exists a solution $u$ of the equation
\begin{equation}\label{eq:browder-minty-surjectivity-equation}
T(u)=g;
\end{equation}
i.e., $T(X)=X^*$.
\end{Theorem}

\begingroup
\renewcommand{\endProof}{\par\nobreak\hfill$\blacksquare$\par\vspace{0.1cm}}
\begin{Proof}
Let $\{x_i\}_{i=1}^\infty$ be a set whose span is dense in $X$. (Recall that we are giving proofs only in the case where $X$ and $X^*$ are separable.) Now let
\begin{equation}\label{eq:browder-minty-dense-finite-dimensional-spaces}
X\supset X_n:=\operatorname{span}\{x_1,x_2,\ldots,x_n\}.
\end{equation}
Our plan is to approximate the problem \eqref{eq:browder-minty-surjectivity-equation} by a finite-dimensional problem
\begin{equation}\label{eq:browder-minty-galerkin-finite-dimensional-problem}
T_n(u_n)=g_n,
\end{equation}
where $u_n\in X_n$, $g_n\in X_n^*$ and $T_n:X_n\to X_n^*$. The idea of approximating infinite-dimensional by finite-dimensional problems is known as \emph{Galerkin\textquotesingle s method}. It is well-known as a device for doing numerical calculations, but is equally useful as a theoretical tool (as we use it here).

We assume $g\in X^*$ is given and define $g_n\in X_n^*$ to be the restriction
\begin{equation}\label{eq:browder-minty-restricted-right-hand-side}
g_n:=g|_{X_n}.
\end{equation}
Similarly, we define the operator $T_n:X_n\to X_n^*$ by
\begin{equation}\label{eq:browder-minty-restricted-nonlinear-mapping}
T_n(u):=T(u)|_{X_n}\quad\text{for all }u\in X_n.
\end{equation}
Another way of saying this is that for any $u,v\in X_n$ we have
\begin{equation}\label{eq:browder-minty-restriction-pairing-identities}
(g-g_n,v)=(T(u)-T_n(u),v)=0.
\end{equation}
The continuity, boundedness and coercivity of $T_n$ are inherited directly from properties of $T$. Thus, using Theorem~\ref{thm:finite-dimensional-coercive-surjectivity}, we see that there exists a sequence $\{u_n\}_{n=1}^\infty$ such that
\begin{equation}\label{eq:browder-minty-galerkin-solution-sequence}
T_n(u_n)=g_n.
\end{equation}
We now use \eqref{eq:browder-minty-restriction-pairing-identities} to note that
\begin{equation}\label{eq:browder-minty-coercive-solution-norm-bound}
\begin{aligned}
\dfrac{(T(u_n),u_n)}{\|u_n\|}&=\dfrac{(T_n(u_n),u_n)}{\|u_n\|}=\dfrac{(g_n,u_n)}{\|u_n\|}\\
&=\dfrac{(g,u_n)}{\|u_n\|}\leq\dfrac{\|g\|\,\|u_n\|}{\|u_n\|}=\|g\|.
\end{aligned}
\end{equation}
Thus, since $T$ is coercive, we must have $\|u_n\|$ bounded. Also, since $T$ is bounded we must have $\|T(u_n)\|$ bounded. Thus, using the weak compactness theorem~\ref{thm:alaoglu-weak-star-subsequence}, the reflexivity of $X$ and the separability of $X^*$, we see that there exists $u\in X$, $\widetilde g\in X^*$ and a subsequence (also labeled $u_n$) such that
\begin{equation}\label{eq:browder-minty-galerkin-weak-solution-limit}
u_n\rightharpoonup u\quad\text{in }X
\end{equation}
and
\begin{equation}\label{eq:browder-minty-galerkin-weak-image-limit}
T(u_n)\rightharpoonup\widetilde g\quad\text{in }X^*.
\end{equation}
We now show that $\widetilde g=g$ and that $T(u)=g$.

We use \eqref{eq:browder-minty-restriction-pairing-identities} again to note that for every basis vector $x_i$,
\begin{equation}\label{eq:browder-minty-identify-weak-image-limit}
(\widetilde g-g,x_i)=\lim_{n\to\infty}(T(u_n)-g,x_i)=0.
\end{equation}
Thus, $\widetilde g=g$.

The key result here is that, once again using \eqref{eq:browder-minty-restriction-pairing-identities}, we can get
\begin{equation}\label{eq:browder-minty-galerkin-weak-pairing-limit}
\begin{aligned}(T(u_n),u_n)&=(T_n(u_n),u_n)\\&=(g_n,u_n)\\&=(g,u_n)\\&\to(g,u).\end{aligned}
\end{equation}
(Since the left-hand side is the product of two \emph{weakly} convergent sequences, this result is nontrivial.) We now use this and the monotonicity of $T$ to get that for every $v\in X$
\begin{equation}\label{eq:browder-minty-monotonicity-limit-inequality}
0\leq(T(u_n)-T(v),u_n-v)\to(T(v)-g,v-u).
\end{equation}
The theorem now follows immediately from Lemma~\ref{lem:banach-minty-equivalence}.
\end{Proof}
\endgroup

\subsubsection{Applications of Monotone Operators to Nonlinear PDEs}\label{subsec:monotone-operators-nonlinear-pde-applications}

In this section we apply the theory of monotone operators to scalar second-order quasilinear PDEs in divergence form. Thus, we assume $\Omega\subset\mathbb{R}^n$ is a domain, and for sufficiently smooth functions $u:\Omega\to\mathbb{R}$, we set
\begin{equation}\label{eq:monotone-quasilinear-divergence-operator}
A(u)=\sum_{|\alpha|\leq1}(-1)^{|\alpha|}D^\alpha\mathcal{A}_\alpha(\bm{x},\delta_1(u(\bm{x}))),
\end{equation}
where
\begin{equation}\label{eq:monotone-quasilinear-first-derivative-jet}
\delta_1(u(\bm{x})):=\{D^\alpha u(\bm{x})\mid|\alpha|\leq1\}.
\end{equation}
We will consider the Dirichlet problem for
\begin{equation}\label{eq:monotone-quasilinear-dirichlet-equation}
A(u)=f
\end{equation}
with $f\in W^{-1,q}(\Omega)$, where $p\in(1,\infty)$ is given and $\frac{1}{p}+\frac{1}{q}=1$. (Here, in analogy to Section~\ref{sec:negative-sobolev-duality}, we define $W^{-1,q}(\Omega)$ to be the dual space of $W_0^{1,p}(\Omega)$.) The ``bivariate form\textquotesingle\textquotesingle\ corresponding to $A$ is
\begin{equation}\label{eq:monotone-quasilinear-bivariate-form}
B(u,v):=\sum_{|\alpha|\leq1}\int_\Omega\mathcal{A}_\alpha(\bm{x},\delta_1(u(\bm{x})))D^\alpha v(\bm{x})\,\mathrm{d}\bm{x}.
\end{equation}


\begin{Definition}
\label{def:monotone-quasilinear-dirichlet-weak-solution}
We say that $u\in W_0^{1,p}(\Omega)$ is a \textbf{weak solution} of the Dirichlet problem for the quasilinear PDE \eqref{eq:monotone-quasilinear-dirichlet-equation} if
\begin{equation}\label{eq:monotone-quasilinear-dirichlet-weak-formulation}
B(u,v)=(f,v)\quad\text{for all }v\in W_0^{1,p}(\Omega).
\end{equation}

\end{Definition}

In order to get an existence theorem we make the following assumptions about the functions $\mathcal{A}_\alpha$.

\begin{enumerate}[label=H-\arabic*.,ref=H-\arabic*]\item\label{hyp:monotone-pde-measurability} For each $|\alpha|\leq1$,
\begin{equation}\label{eq:monotone-pde-coefficient-measurability}
\bm{x}\mapsto\mathcal{A}_\alpha(\bm{x},\delta_1)
\end{equation}
is measurable for every fixed $\delta_1=(\bm{x}i_\alpha)_{|\alpha|\leq1}\in\mathbb{R}^{n+1}$.

\item\label{hyp:monotone-pde-continuity} For each $|\alpha|\leq1$,
\begin{equation}\label{eq:monotone-pde-coefficient-jet-continuity}
\delta_1\mapsto\mathcal{A}_\alpha(\bm{x},\delta_1)
\end{equation}
is in $C(\mathbb{R}^{n+1})$ for almost every $\bm{x}\in\Omega$.

\item\label{hyp:monotone-pde-monotonicity} For every $\delta_1^1=(\bm{x}i_\alpha^1)_{|\alpha|\leq1}\in\mathbb{R}^{n+1}$, $\delta_1^2=(\bm{x}i_\alpha^2)_{|\alpha|\leq1}\in\mathbb{R}^{n+1}$ and every $\bm{x}\in\Omega$, we have
\begin{equation}\label{eq:monotone-pde-coefficient-monotonicity}
\sum_{|\alpha|\leq1}[\mathcal{A}_\alpha(\bm{x},\delta_1^1)-\mathcal{A}_\alpha(\bm{x},\delta_1^2)](\bm{x}i_\alpha^1-\bm{x}i_\alpha^2)\geq0.
\end{equation}


\item\label{hyp:monotone-pde-coercivity} There exists $p\in(1,\infty)$, a constant $c_0>0$, a function $h\in L^1(\Omega)$ such that for every $\bm{x}\in\Omega$ and every $\delta_1=(\bm{x}i_\alpha)_{|\alpha|\leq1}\in\mathbb{R}^{n+1}$ we have
\begin{equation}\label{eq:monotone-pde-coefficient-coercivity}
\sum_{|\alpha|\leq1}\mathcal{A}_\alpha(\bm{x},\delta_1)\bm{x}i_\alpha\geq c_0|\delta_1|^p-h(\bm{x}).
\end{equation}


\item\label{hyp:monotone-pde-growth} There exists a constant $c_1>0$ and a function $g\in L^q(\Omega)$, ($q:=p/(p-1)$) such that for every $\bm{x}\in\Omega$ and every $\delta_1=(\bm{x}i_\alpha)_{|\alpha|\leq1}\in\mathbb{R}^{n+1}$ we have
\begin{equation}\label{eq:monotone-pde-coefficient-growth-bound}
|\mathcal{A}_\alpha(\bm{x},\delta_1)|\leq c_1|\delta_1|^{p-1}+g(\bm{x}).
\end{equation}
\end{enumerate}

Under these assumptions we get the following result.

\begin{Theorem}
\label{thm:monotone-quasilinear-dirichlet-existence}
Let $\mathcal{A}_\alpha$ satisfy hypotheses \ref{hyp:monotone-pde-measurability}--\ref{hyp:monotone-pde-growth}, and let $p$ and $q$ be as defined in hypotheses \ref{hyp:monotone-pde-coercivity} and \ref{hyp:monotone-pde-growth}, respectively. Then, for every $f\in W^{-1,q}(\Omega)$, there exists a weak solution $u\in W_0^{1,p}(\Omega)$ of the Dirichlet problem for the quasilinear PDE \eqref{eq:monotone-quasilinear-dirichlet-equation}.
\end{Theorem}

We break the proof of this theorem into a series of lemmas, the first of which is the following.

\begin{Lemma}
\label{lem:monotone-pde-induced-bounded-dual-mapping}
For each fixed $u\in W_0^{1,p}(\Omega)$, the mapping
\begin{equation}\label{eq:monotone-pde-bounded-test-functional}
W_0^{1,p}(\Omega)\ni v\mapsto B(u,v)\in\mathbb{R}
\end{equation}
is a bounded linear functional. Thus, there exists a mapping
\begin{equation}\label{eq:monotone-pde-induced-dual-mapping}
X=W_0^{1,p}(\Omega)\ni u\mapsto T(u)\in W^{-1,q}(\Omega)=X^*
\end{equation}
such that
\begin{equation}\label{eq:monotone-pde-induced-duality-pairing}
B(u,v)=(T(u),v),
\end{equation}
for all $u,v$ in $W_0^{1,p}(\Omega)$. Furthermore, the nonlinear mapping is bounded.
\end{Lemma}

\begingroup
\renewcommand{\endProof}{\par\nobreak\hfill$\blacksquare$\par\vspace{0.1cm}}
\begin{Proof}
The linearity of the mapping $v\mapsto B(u,v)$ is obvious. If $u,v\in W_0^{1,p}(\Omega)$, then, using \ref{hyp:monotone-pde-growth} and Hölder\textquotesingle s inequality, we get the estimate:\[\begin{aligned}
|B(u,v)|&\leq\sum_{|\alpha|\leq1}\int_\Omega|\mathcal{A}_\alpha(\bm{x},\delta_1(u(\bm{x})))D^\alpha v(\bm{x})|\,\mathrm{d}\bm{x}\\
&\leq\sum_{|\alpha|\leq1}\int_\Omega[c_1|\delta_1(u(\bm{x}))|^{p-1}+g(\bm{x})]|D^\alpha v(\bm{x})|\,\mathrm{d}\bm{x}\\
&\leq\sum_{|\alpha|\leq1}c_1\left\{\int_\Omega|\delta_1(u(\bm{x}))|^{q(p-1)}\right\}^{1/q}\left\{\int_\Omega|D^\alpha v|^p\,\mathrm{d}\bm{x}\right\}^{1/p}\\
&\quad+\sum_{|\alpha|\leq1}\left\{\int_\Omega|g|^q\,\mathrm{d}\bm{x}\right\}^{1/q}\left\{\int_\Omega|D^\alpha v|^p\,\mathrm{d}\bm{x}\right\}^{1/p}\\
&\leq C(\|u\|_{1,p}^{p-1}+1)\|v\|_{1,p}.
\end{aligned}\]This shows that for each $u\in W_0^{1,p}(\Omega)$ the mapping $v\mapsto B(u,v)$ is bounded. This implies the existence of the mapping $T$.

Furthermore, the same estimate shows that $T$ is bounded since
\begin{equation}\label{eq:monotone-pde-induced-dual-norm-bound}
\|T(u)\|_{-1,q}=\sup_{\substack{v\in W_0^{1,p}(\Omega)\\\|v\|_{1,p}=1}}|(T(u),v)|\leq C(\|u\|_{1,p}^{p-1}+1).
\end{equation}
Thus, $T$ maps bounded sets in $W_0^{1,p}(\Omega)$ into bounded sets in $W^{-1,q}(\Omega)$. This completes the proof.
\end{Proof}
\endgroup

\begin{Lemma}
\label{lem:monotone-pde-induced-mapping-monotone-coercive}
The operator $T$ defined in Lemma~\ref{lem:monotone-pde-induced-bounded-dual-mapping} is monotone and coercive.
\end{Lemma}

\begingroup
\renewcommand{\endProof}{\par\nobreak\hfill$\blacksquare$\par\vspace{0.1cm}}
\begin{Proof}
To prove monotonicity we use hypothesis \ref{hyp:monotone-pde-monotonicity} to get\[\begin{aligned}
(T(u)-T(v),u-v)&=B(u,u-v)-B(v,u-v)\\
&=\sum_{|\alpha|\leq1}\int_\Omega[\mathcal{A}_\alpha(\bm{x},\delta_1(u))-\mathcal{A}_\alpha(\bm{x},\delta_1(v))][D^\alpha u-D^\alpha v]\,\mathrm{d}\bm{x}\\
&\geq0.
\end{aligned}\]To prove coercivity we use hypothesis \ref{hyp:monotone-pde-coercivity} to get\[\begin{aligned}
\dfrac{(T(u),u)}{\|u\|_{1,p}}&=\dfrac{B(u,u)}{\|u\|_{1,p}}\\
&=\|u\|_{1,p}^{-1}\sum_{|\alpha|\leq1}\int_\Omega\mathcal{A}_\alpha(\bm{x},\delta_1(u))D^\alpha u\,\mathrm{d}\bm{x}\\
&\geq\|u\|_{1,p}^{-1}\left[c_0\int_\Omega|\delta_1(u)|^p\,\mathrm{d}\bm{x}-\int_\Omega h(\bm{x})\,\mathrm{d}\bm{x}\right]\\
&\geq\|u\|_{1,p}^{-1}C[\|u\|_{1,p}^p-1].
\end{aligned}\]Since $p>1$ we have
\begin{equation}\label{eq:monotone-pde-induced-mapping-coercivity-limit}
\dfrac{(T(u),u)}{\|u\|_{1,p}}\to\infty\quad\text{as }\|u\|_{1,p}\to\infty.
\end{equation}
This completes the proof.
\end{Proof}
\endgroup

Thus, to apply the Browder-Minty theorem to the mapping $T$ and complete the proof of Theorem~\ref{thm:monotone-quasilinear-dirichlet-existence} we need only show the following.

\begin{Lemma}
\label{lem:monotone-pde-induced-mapping-continuity}
The mapping $T:W_0^{1,p}(\Omega)\to W^{-1,q}(\Omega)$ is continuous.
\end{Lemma}

In the next section we describe a tool called \emph{Nemytskii operators} which we can use to prove this lemma.

\begin{Example}
\label{ex:anisotropic-p-power-nonlinear-elliptic-operator}
Consider the second-order nonlinear partial differential operator
\begin{equation}\label{eq:anisotropic-p-power-divergence-example}
A(u):=-\sum_{i=1}^n\dfrac{\partial}{\partial x_i}\left(\left|\dfrac{\partial u}{\partial x_i}\right|^{p-2}\dfrac{\partial u}{\partial x_i}\right)+|u|^{p-2}u,
\end{equation}
where $p\in(1,\infty)$. Note that the case $p=2$ is simply the Laplacian plus a lower-order term which we have already considered in our material on linear problems. Here,\[\begin{aligned}
\mathcal{A}_{(0,0,\ldots,0)}&=a_0(\bm{x},u,u_{x_1},\ldots,u_{x_n})=|u|^{p-2}u,\\
\mathcal{A}_{(0,\ldots,0,1,0,\ldots,0)}&=a_i(\bm{x},u,u_{x_1},\ldots,u_{x_n})=|u_{x_i}|^{p-2}u_{x_i},\quad i=1,\ldots,n.
\end{aligned}\]We wish to verify that these $\mathcal{A}_\alpha$ satisfy the hypotheses \ref{hyp:monotone-pde-measurability} to \ref{hyp:monotone-pde-growth}. Hypotheses \ref{hyp:monotone-pde-measurability} and \ref{hyp:monotone-pde-continuity} obviously hold. To verify \ref{hyp:monotone-pde-monotonicity} we let $\delta_1^1=(\bm{x}i_0^1,\bm{x}i_1^1,\ldots,\bm{x}i_n^1)$ and $\delta_1^2=(\bm{x}i_0^2,\bm{x}i_1^2,\ldots,\bm{x}i_n^2)$ and calculate\[\begin{aligned}
&\sum_{|\alpha|\leq1}[\mathcal{A}_\alpha(\bm{x},\delta_1^1)-\mathcal{A}_\alpha(\bm{x},\delta_1^2)](\bm{x}i_\alpha^1-\bm{x}i_\alpha^2)\\
&\quad=\sum_{i=0}^n\bigl(|\bm{x}i_i^1|^{p-2}\bm{x}i_i^1-|\bm{x}i_i^2|^{p-2}\bm{x}i_i^2\bigr)(\bm{x}i_i^1-\bm{x}i_i^2)\\
&\quad\geq0.
\end{aligned}\]To verify \ref{hyp:monotone-pde-coercivity} we let $\delta_1=(\bm{x}i_0,\bm{x}i_1,\ldots,\bm{x}i_n)$ and get
\begin{equation}\label{eq:anisotropic-p-power-printed-coercivity-chain}
\sum_{|\alpha|\leq1}\mathcal{A}_\alpha(\bm{x},\delta_1)\bm{x}i_\alpha=\sum_{i=0}^n|\bm{x}i_i|^{p-2}\bm{x}i_i\bm{x}i_i=\sum_{i=1}^n|\bm{x}i_i|^p\geq c_0|\delta_i|^p
\end{equation}
We see that \ref{hyp:monotone-pde-growth} holds since
\begin{equation}\label{eq:anisotropic-p-power-coefficient-growth}
|\mathcal{A}_\alpha(\bm{x},\delta_1)|=|a_i(\bm{x},\delta_1)|=|\bm{x}i_i|^{p-1}\leq|\delta_1|^{p-1}.
\end{equation}
Thus the following existence result follows immediately from Theorem~\ref{thm:monotone-quasilinear-dirichlet-existence}.
\end{Example}

\begin{Theorem}
\label{thm:anisotropic-p-power-dirichlet-existence}
Let the nonlinear second order partial differential operator $A$ be defined by \eqref{eq:anisotropic-p-power-divergence-example}. Then for every $f\in W^{-1,q}(\Omega)$ there exists a weak solution $u\in W_0^{1,p}$ of the equation
\begin{equation}\label{eq:anisotropic-p-power-dirichlet-existence-equation}
A(u)=f.
\end{equation}

\end{Theorem}

\subsubsection{Nemytskii Operators}\label{subsec:nemytskii-operators}

In the following section we state without proof some important results on the composition of $L^p(\Omega)$ with nonlinear functions. For a more detailed treatment, the reader could consult \cite{Li}.

\begin{Definition}
\label{def:caratheodory-conditions}
Let $\Omega\subset\mathbb{R}^n$ be a domain. We say that a function
\begin{equation}\label{eq:caratheodory-nonlinear-function-domain}
\Omega\times\mathbb{R}^m\ni(\bm{x},\bm{u})\mapsto f(\bm{x},\bm{u})\in\mathbb{R}
\end{equation}
satisfies the \textbf{Carathéodory conditions} if
\begin{equation}\label{eq:caratheodory-value-variable-continuity}
\bm{u}\mapsto f(\bm{x},\bm{u})\text{ is continuous for almost every }\bm{x}\in\Omega
\end{equation}
and
\begin{equation}\label{eq:caratheodory-spatial-measurability-printed-domain}
\bm{x}\mapsto f(\bm{x},\bm{u})\text{ is measurable for every }\bm{u}\in\Omega.
\end{equation}

\end{Definition}

Given any $f$ satisfying the Carathéodory conditions and a function $\bm{u}:\Omega\to\mathbb{R}^m$, we can define another function by composition
\begin{equation}\label{eq:nemytskii-pointwise-composition-operator}
\mathcal{F}(\bm{u})(\bm{x}):=f(\bm{x},\bm{u}(\bm{x})).
\end{equation}
The composition operator $\mathcal{F}$ is called a \emph{Nemytskii operator}. Our main theorem is on the boundedness and continuity of these operators from $L^p(\Omega)$ to $L^q(\Omega)$.

\begin{Theorem}
\label{thm:nemytskii-lp-lq-bounded-continuous}
Let $\Omega\subset\mathbb{R}^n$ be a domain, and let
\begin{equation}\label{eq:nemytskii-growth-theorem-function-domain}
\Omega\times\mathbb{R}^m\ni(\bm{x},\bm{u})\mapsto f(\bm{x},\bm{u})\in\mathbb{R}
\end{equation}
satisfy the Carathéodory conditions. In addition, let $p\in(1,\infty)$ and $g\in L^q(\Omega)$ (where $\frac{1}{p}+\frac{1}{q}=1$) be given, and let $f$ satisfy
\begin{equation}\label{eq:nemytskii-growth-bound-lp-lq}
|f(\bm{x},\bm{u})|\leq C|\bm{u}|^{p-1}+g(\bm{x}).
\end{equation}
Then the Nemytskii operator $\mathcal{F}$ defined by \eqref{eq:nemytskii-pointwise-composition-operator} is a bounded and continuous map from $L^p(\Omega)$ to $L^q(\Omega)$.
\end{Theorem}

\begin{Remark}
\label{rem:nemytskii-monotone-pde-continuity-corollary}
Lemma~\ref{lem:monotone-pde-induced-mapping-continuity} follows as a corollary to this theorem. To see this we simply need to apply hypotheses \ref{hyp:monotone-pde-measurability}, \ref{hyp:monotone-pde-continuity} and \ref{hyp:monotone-pde-growth} to see that each $\mathcal{A}_\alpha$ can be used as a Nemytskii operator satisfying the appropriate growth conditions. The continuity of $T$ from $W^{1,p}(\Omega)$ to $W^{-1,q}(\Omega)$ follows from the continuity of $\delta_1(\bm{x})\mapsto\mathcal{A}_\alpha(\bm{x},\delta_1(\bm{x}))$ as a map from $L^p(\Omega)$ to $L^q(\Omega)$.
\end{Remark}
